\documentclass[11pt]{article}

\usepackage{maiacustom}

\begin{document}

\psettitle{Astronomy Question Bank}

These questions were carefully written or selected to prepare students for the astronomy team selection process in Brazil. Some questions are not original; those are marked with () before the statement. The question-bank template is the same one used by Professor Kevin Zhou. His work is invaluable, and many ideas in this set can be found in his handouts.

% \begin{psidea}{Idea title}{}
% Idea
% \end{psidea}

% \begin{psexample}{Example title}{}
% Example
% \end{psexample}

% \begin{pssolution*}{}{}
% Solution
% \end{pssolution*}

% \begin{psremark*}{Remark title}{}
% Remark
% \end{psremark*}
\section{Miscellaneous}
\pts{4}
\begin{pproblem} \clock{1}{0}
    Dudu Leiteiro was observing the sky in the interior of his farm in Mato Grosso do Sul and observed the binary system formed by the stars Iaum and Sezenem. Dudu observed the stars and collected the following data, with an interval of 6 months between them:
    \\
    \begin{center}
        \begin{tabular}{|c|c|c|}
            \hline % Top horizontal line
            Measurement & Iaum & Sezenem   \\ 
            \hline
            Right ascension \(1\) & \(4^h19^m53.91^s\) & \(4^h19^m53.078^s\) \\
            Right ascension \(2\) & \(4^h19^m53.92^s\) & \(4^h19^m53.077^s\) \\
            Declination \(1\) & -\(13^\circ \ 33' \ 45.28''\) & -\(13^\circ \ 33' \ 47.78''\) \\
            Declination \(2\) & -\(13^\circ \ 33' \ 45.20''\) & -\(13^\circ \ 33' \ 48.03''\) \\
            \hline
        \end{tabular} 
\end{center}
    With these data, you and Dudu Leiteiro together will analyze properties of this binary system. The first important step is to find the coordinates of the center of mass of the system, denoted by the subscript \(_{CM}\). You remember that, in one of the classes on the study of binaries, your professor, LuCav, taught you that
    \[
    \delta_{CM} = \delta_A + \Delta\delta_A
    \]

    \[
    \alpha_{CM} = \alpha_A + \Delta\alpha_A
    \]

    where \(\delta_A\) and \(\alpha_A\) are the coordinates of one of the stars in the binary and \(\Delta\) represents the difference in coordinates between the star and the center of mass. We have the following relation:

    \[
    \frac{a_A}{a} = \frac{\Delta\delta_A}{\Delta\delta_{CM}} = \frac{\Delta\alpha_A}{\Delta\alpha_{CM}}
    \]

    where \(a_A\) is the distance from star \(A\) to the \(CM\), \(a\) is the semi-major axis of the orbit, and \(\Delta x_{CM}\) is the variation of the coordinates of the \(CM\). 
    
    \textbf{a)} Given that Iaum has \(3/2\) the mass of Sezenem, calculate \(\Delta\delta_{CM}\) and \(\Delta\alpha_{CM}\).
    
    \textbf{b)} With this, and considering that the angular variations are small enough for spherical triangles to be treated as plane, find the parallax of the system and its distance to Earth.

    \textbf{c)} Considering the mass of Iaum, \(M_I = 4.9 M_\odot\), and the period of the system equal to \(P = 29.01\) years, calculate the largest redshift coming from Sezenem, given that both orbits are circular and Sezenem has a tangential velocity of \(\mu = 1509''/\text{year}\).

\end{pproblem}

\pts{3}
\begin{pproblem} Marisso was tired of not being able to find the position of a star precisely with his telescope and decided to investigate the effects that could be causing this apparent error. After reading some papers, he discovered \(2\) main effects that make an object appear to lie at an angle \(\Delta \theta_i\), offset from its original position. They are \textit{parallax} and \textit{stellar aberration}. In this problem, your goal is to help Marrisso understand why these effects occur.
    \begin{alternativas}
        \item Parallax is the most basic of them, and it occurs because of the change in the position of the Earth over the course of the year. Consider a star located so that the line Sun--star is perpendicular to the plane of the Earth's orbit. Draw a diagram of the situation and, taking the radius of the Earth's orbit as \(r\) and the distance to the star as \(d\), obtain a formula for \(\Delta \theta_p\) caused by parallax.
        
        \item Now suppose that the line Sun--star makes an arbitrary angle \(\phi\) with the Earth's orbit. How does your answer change?
        
        Stellar aberration, in turn, comes from relativistic effects, which are explored below.

        \item Consider a frame \(S'\) moving with velocity \(v\hat{\mathbf{x}}\) relative to the frame \(S\). How are the coordinates \((x', t')\) related to the coordinates \((x,t)\)? Leave your answers in terms of \(\gamma\).
        
        \item Suppose there is a radiation emitter in the frame \(S'\) and that it emits light at an angle \(\alpha'\) relative to the \(x\) axis. In the frame \(S\), the device will appear to emit light at an angle \(\alpha\). Prove, using the Lorentz transformations, that the relation between \(\alpha\) and \(\alpha'\) is given by
        
        \[\cos\alpha = \frac{\cos\alpha' + v/c}{1+(\cos\alpha') v/c}\]

        \item Repeat the previous item, but prove it using relativistic velocity addition.

        \item Considering that the line Sun--star is perpendicular to the plane of the Earth's orbit, and that the Earth moves with velocity \(v\), find an expression for the shift \(\Delta\theta_A\) caused by stellar aberration.
    
        \item Which of these effects do you think is more significant, parallax or stellar aberration?
    \end{alternativas}

    \begin{pssolution*}{}{}
        \begin{alternativas}
            \item % too lazy to fetch GeoGebra
            
            \item % lazy GeoGebra
                
            \item The Lorentz transformations (using \(c=1\) ) are given by
            \[x = \gamma (x' + vt'), \ \ t = \gamma(t' + vx')\]

            \item Note that in the frame \(S\), we have \(\cos \alpha = \frac{x}{t}\). In the frame \(S'\) we have
            
            \[x' = t'\cos\alpha'\]

            Using the Lorentz transformations,

            \[t = \gamma t'(1+v\cos\alpha') \]

            \[x = \gamma t'(\cos\alpha'+ v)\]

            Using \(\cos\alpha = x/t\), we have

            \[\cos\alpha = \frac{\cos\alpha' + v}{1+ v\cos\alpha'}\]

            ``Putting the \(c\)'s back'', we have

            \[\boxed{\cos\alpha = \frac{\cos\alpha' + v/c}{1+(v/c)\cos\alpha'}}\]


            \item In the frame \(S'\), the velocity along \(x'\) is given by \(v_x' = c\cos\alpha'\). In the frame \(S\), we have
            
            \[v_x = "v_x'+v" = \frac{v_x'+v}{1+(v_x'v)/c^2}\]

            Substituting \(v_x'\) and using \(\cos\alpha = v_x/c\), we have

            \[\boxed{\cos\alpha =\frac{\cos\alpha' + v/c}{1+(v/c)\cos\alpha'}} \]

            \item In the Earth's frame (which is moving and therefore corresponds to the frame \(S'\)), we have \(\alpha' = \pi/2\). Plugging this into the answer, we obtain
            
            \[\cos\alpha = v/c \rightarrow \sin\alpha = \sqrt{1-(v/c)^2} = \frac{1}{\gamma}\]

            Using the small-angle approximation, \(\boxed{\alpha \approx 1/\gamma}\).

            \item While the Sun--Earth distance is on the order of a few light-minutes, the smallest distance between the Earth and another star is on the order of light-years, so \(\theta \sim 10^{-6}\). The speed of the Earth is approximately \(30\) km/s, so \(v/c \sim 10^{-4}\). Therefore the effects of aberration are more apparent than those of parallax.


        \end{alternativas}

        \end{pssolution*}

\end{pproblem}
\end{document}
